\documentclass[aspectratio=169,10pt]{beamer}
\makeatletter\def\input@path{{../../estilo/}{../estilo/}{estilo/}}\makeatother
\usepackage{devmobile}
\modulo[Kotlin]{08}{Arquitetura de Apps}{Camadas, MVVM, UiState, repositórios, casos de uso, injeção de dependência com Hilt e organização do projeto}

\begin{document}
\capa

\begin{frame}{Objetivos do módulo}
  \begin{conceito}[Ao final deste módulo você será capaz de]
    \begin{itemize}
      \item Explicar os princípios da arquitetura recomendada para apps Android.
      \item Separar o app em camadas de UI, domínio e dados.
      \item Implementar o padrão MVVM com \cod{ViewModel}, \cod{UiState} e eventos.
      \item Criar repositórios como fonte única da verdade.
      \item Encapsular regras de negócio em casos de uso.
      \item Configurar injeção de dependência com Hilt e organizar pacotes e módulos.
    \end{itemize}
  \end{conceito}
\end{frame}

\sumario

% =============================================================================
%  Fundamentos compartilhados --- Módulo 08: arquitetura de software
%  Incluído pelas trilhas Kotlin e Flutter com % =============================================================================
%  Fundamentos compartilhados --- Módulo 08: arquitetura de software
%  Incluído pelas trilhas Kotlin e Flutter com % =============================================================================
%  Fundamentos compartilhados --- Módulo 08: arquitetura de software
%  Incluído pelas trilhas Kotlin e Flutter com \input{../fundamentos/f08-arquitetura}
%  (coesão, acoplamento, SOLID e padrões de projeto estão no Módulo 03)
% =============================================================================
\section{Fundamentos: arquitetura de software}

\begin{frame}{O que é arquitetura de software}
  \begin{columns}[T]
    \begin{column}{0.52\textwidth}
      \begin{definicao}[Arquitetura]
        As grandes estruturas de um sistema: seus elementos, as relações entre eles e as decisões \textbf{difíceis de mudar depois}. Ela existe para proteger atributos de qualidade.
      \end{definicao}
      \scriptsize
      \begin{tabularx}{\linewidth}{@{}l L@{}}
        \toprule
        \cabecalho{Atributo} & \cabecalho{Pergunta no LabMonitor} \\
        \midrule
        Manutenibilidade & trocar a API exige mexer nas telas? \\
        Testabilidade & dá para testar a regra de reserva sem um aparelho? \\
        Desempenho & a lista de leituras rola sem engasgar? \\
        Funcionamento \emph{offline} & o app mostra dados sem internet? \\
        Segurança & o \emph{token} fica protegido? \\
        \bottomrule
      \end{tabularx}
    \end{column}
    \begin{column}{0.46\textwidth}
      \begin{analogia}
        A planta de uma casa decide por onde passam os canos antes da primeira parede. Pintar um quarto não é arquitetura; mudar a cozinha de lugar é.
      \end{analogia}
      \begin{conceito}[Boa arquitetura em um app]
        Permite mudar a interface sem tocar nos dados, trocar a API sem tocar na interface e testar regras sem abrir o app.
      \end{conceito}
      \begin{dica}
        Coesão, acoplamento, SOLID e padrões de projeto foram vistos no Módulo 03: aqui eles ganham escala.
      \end{dica}
    \end{column}
  \end{columns}
\end{frame}

\begin{frame}{Padrões de apresentação: de MVC a MVI}
  \centering
  \begin{tikzpicture}[every node/.style={font=\tiny}]
    \draw[seta] (0,0) -- (13.4,0);
    \foreach \x/\t/\d in {1.2/MVC/{1979, Smalltalk\\ (Trygve Reenskaug)},4.6/MVP/{anos 1990 (Taligent);\\ comum no Android até 2017},8.1/MVVM/{2005, Microsoft (WPF);\\ padrão no Android e no Flutter},11.6/MVI/{anos 2010, inspirado\\ no Redux e no Cycle.js}}{
      \fill[ambar] (\x,0) circle (2pt);
      \node[font=\scriptsize\bfseries,text=azulpetroleo,above=1mm] at (\x,0) {\t};
      \node[rotulo,below=1mm] at (\x,0) {\d};
    }
  \end{tikzpicture}
  \raggedright
  \scriptsize
  \begin{tabularx}{\textwidth}{@{}l L L L@{}}
    \toprule
    \cabecalho{Padrão} & \cabecalho{Quem atualiza a tela} & \cabecalho{Ponto forte} & \cabecalho{Ponto fraco} \\
    \midrule
    MVC & o \emph{controller} manipula a \emph{view} & simples & no Android, a \cod{Activity} virava um \emph{controller} gigante \\
    MVP & o \emph{presenter} chama métodos de uma interface da \emph{view} & testável sem a tela & muito código de interface; guarda referência à \emph{view} \\
    MVVM & a \emph{view} observa o estado exposto pelo \emph{view model} & casa com UI declarativa e sobrevive à rotação & vários fluxos de estado podem se desencontrar \\
    MVI & estado único e imutável; cada intenção gera um novo estado & previsível e fácil de depurar & cerimônia demais para telas simples \\
    \bottomrule
  \end{tabularx}
  \normalsize
  \begin{conceito}[Na prática atual]
    Compose e Flutter usam MVVM com um estado único por tela (\cod{UiState}): uma mistura de MVVM com a disciplina do MVI.
  \end{conceito}
\end{frame}

\begin{frame}{Arquitetura limpa e arquitetura hexagonal}
  \begin{columns}[c]
    \begin{column}{0.46\textwidth}
      \centering
      \begin{tikzpicture}[every node/.style={font=\tiny,align=center}]
        \fill[cinzaclaro] (0,0) circle (2.15);
        \fill[cinzafundo,draw=azulpetroleo!40] (0,0) circle (1.6);
        \fill[ambarclaro!60,draw=ambar] (0,0) circle (1.05);
        \fill[azulpetroleo] (0,0) circle (0.55);
        \node[text=white,font=\tiny\bfseries] at (0,0) {Entidades};
        \node at (0,0.8) {Casos de uso};
        \node at (0,1.33) {Adaptadores};
        \node at (0,1.88) {Frameworks e drivers};
        \node[rotulo] at (0,-1.33) {\emph{view models}, repositórios};
        \node[rotulo] at (0,-1.88) {UI, banco, HTTP, BLE};
        \draw[setaambar] (2.6,-1.2) -- (1.2,-0.55);
        \node[rotulo,text width=1.5cm] at (3.05,-1.5) {dependências apontam para dentro};
      \end{tikzpicture}
    \end{column}
    \begin{column}{0.52\textwidth}
      \begin{definicao}[Regra da dependência]
        O código das camadas internas não conhece as externas: as regras de negócio não importam nada de Android, Flutter, HTTP ou SQL.
      \end{definicao}
      \begin{historia}
        Alistair Cockburn propôs a arquitetura hexagonal, de portas e adaptadores, em 2005. Robert C. Martin reuniu ideias parecidas na Arquitetura Limpa, em 2012.
      \end{historia}
      \begin{dica}[Na medida certa]
        Em apps pequenos, camadas de UI e de dados bastam. Crie a camada de domínio quando regras se repetirem ou precisarem de testes isolados, como recomendam o Android e o Flutter.
      \end{dica}
    \end{column}
  \end{columns}
\end{frame}

\begin{frame}[fragile]{Inversão de dependência na prática}
  \begin{columns}[T]
    \begin{column}{0.5\textwidth}
      \centering
      \begin{tikzpicture}[every node/.style={font=\tiny}]
        \node[rotulo,anchor=west] at (-1.1,0.55) {\textbf{sem inversão}};
        \node[caixaescura,font=\tiny,minimum width=2cm] (r1) at (0,0) {regras de leitura};
        \node[caixa,font=\tiny,minimum width=1.8cm] (b1) at (3.4,0) {\cod{FonteBle}};
        \draw[seta] (r1) -- node[rotulo,above]{conhece} (b1);
        \node[rotulo,anchor=west] at (-1.1,-0.75) {\textbf{com inversão}};
        \node[caixaescura,font=\tiny,minimum width=2cm] (r2) at (0,-1.3) {regras de leitura};
        \node[caixaambar,font=\tiny,minimum width=2.4cm] (i2) at (3.4,-1.3) {\emph{interface}\\ \cod{FonteLeituras}};
        \node[caixa,font=\tiny,minimum width=1.35cm] (b2) at (1.95,-2.75) {\cod{FonteBle}};
        \node[caixa,font=\tiny,minimum width=1.35cm] (m2) at (3.4,-2.75) {\cod{FonteMqtt}};
        \node[caixa,font=\tiny,minimum width=1.35cm] (f2) at (4.85,-2.75) {\cod{FonteFake}};
        \draw[seta] (r2) -- node[rotulo,above]{usa} (i2);
        \foreach \n in {b2,m2,f2}{\draw[setaambar] (\n.north) -- (\n.north |- i2.south);}
        \node[rotulo,below=1mm of m2] {implementam a interface};
      \end{tikzpicture}
\begin{codigo}[title=Dart]{Dart}
abstract interface class FonteLeituras {
  Stream<Leitura> leituras();
}
\end{codigo}
    \end{column}
    \begin{column}{0.48\textwidth}
      \begin{definicao}[Inversão de dependência]
        Módulos de alto nível não dependem de detalhes; ambos dependem de abstrações. Quem precisa define a interface; a camada de dados a implementa.
      \end{definicao}
\begin{codigo}[title=Kotlin]{Kotlin}
interface FonteLeituras {
  fun leituras(): Flow<Leitura>
}
\end{codigo}
      \begin{conceito}[O que se ganha]
        Trocar BLE por MQTT sem mexer nas regras, testar com uma fonte falsa e fazer a tela antes de o hardware existir.
      \end{conceito}
    \end{column}
  \end{columns}
\end{frame}

\begin{frame}{Injeção de dependência e localizador de serviços}
  \begin{columns}[T]
    \begin{column}{0.5\textwidth}
      \begin{definicao}[Inversão de controle]
        A classe não cria o que usa: alguém de fora monta os objetos e os entrega. É o princípio de Hollywood: ``não nos chame, nós chamamos você''.
      \end{definicao}
      \centering
      \begin{tikzpicture}[every node/.style={font=\tiny}]
        \node[caixaambarforte,font=\tiny\bfseries,minimum width=5.6cm] (raiz) at (0,0) {raiz de composição (\cod{Application}, \cod{main})};
        \node[caixa,font=\tiny,minimum width=1.7cm] (h) at (-1.9,-1) {cliente HTTP};
        \node[caixa,font=\tiny,minimum width=1.7cm] (b) at (0,-1) {banco};
        \node[caixapetroleo,font=\tiny,minimum width=1.9cm] (r) at (-0.95,-2) {repositório};
        \node[caixaescura,font=\tiny,minimum width=1.7cm] (v) at (1.9,-2) {\emph{view model}};
        \draw[setatrac] (raiz.south -| h) -- (h); \draw[setatrac] (raiz.south -| b) -- (b);
        \draw[seta] (h) -- (r); \draw[seta] (b) -- (r);
        \draw[seta] (r) -- node[rotulo,above]{injeta} (v);
        \node[rotulo] at (1.9,-1.05) {tracejado: cria};
      \end{tikzpicture}
      \raggedright
      \begin{historia}
        Martin Fowler batizou o padrão de injeção de dependência em 2004.
      \end{historia}
    \end{column}
    \begin{column}{0.48\textwidth}
      \scriptsize
      \begin{tabularx}{\linewidth}{@{}l L@{}}
        \toprule
        \cabecalho{Abordagem} & \cabecalho{Como funciona} \\
        \midrule
        Construtor & recebe tudo pronto; dependências explícitas e fáceis de trocar em testes \\
        Contêiner & uma biblioteca monta o grafo: Hilt, Koin; Provider, Riverpod \\
        Localizador & a classe pede a um registro global (\cod{get\_it}): prático, mas esconde dependências e só falha em execução \\
        \bottomrule
      \end{tabularx}
      \normalsize
      \begin{pordentro}[Tempo de vida (escopo)]
        \textbf{Único no app}: cliente HTTP, banco. \textbf{Por tela}: \emph{view model}. \textbf{Por uso}: objeto novo a cada pedido. Escopo errado causa vazamento ou estado compartilhado por engano.
      \end{pordentro}
    \end{column}
  \end{columns}
\end{frame}

\begin{frame}{Decisões de arquitetura: trocas e registros}
  \begin{columns}[T]
    \begin{column}{0.5\textwidth}
      \begin{conceito}[Toda decisão é uma troca]
        Mais camadas trazem isolamento e testabilidade, mas também mais arquivos e indireção. Pese o tamanho do time, a vida esperada do app e o custo de errar.
      \end{conceito}
      \scriptsize
      \begin{tabularx}{\linewidth}{@{}l L@{}}
        \toprule
        \cabecalho{Princípio} & \cabecalho{Significado} \\
        \midrule
        YAGNI & não construa o que ainda não é necessário \\
        KISS & prefira a solução mais simples que funciona \\
        DRY & cada conhecimento em um só lugar \\
        Regra de três & abstraia na terceira repetição, não na primeira \\
        \bottomrule
      \end{tabularx}
    \end{column}
    \begin{column}{0.48\textwidth}
      \begin{definicao}[ADR (\emph{Architecture Decision Record})]
        Texto curto, versionado com o código, que registra uma decisão: contexto, opções, decisão e consequências. Responde no futuro: ``por que fizemos assim?''.
      \end{definicao}
      \begin{exemplo}[ADR 007 --- Estado das telas]
        \textbf{Contexto}: leituras do sensor em tempo real. \textbf{Decisão}: um \cod{UiState} imutável por tela, exposto como fluxo observável. \textbf{Consequência}: telas testáveis sem interface.
      \end{exemplo}
      \begin{historia}
        Michael Nygard propôs os ADRs em 2011.
      \end{historia}
    \end{column}
  \end{columns}
\end{frame}

%  (coesão, acoplamento, SOLID e padrões de projeto estão no Módulo 03)
% =============================================================================
\section{Fundamentos: arquitetura de software}

\begin{frame}{O que é arquitetura de software}
  \begin{columns}[T]
    \begin{column}{0.52\textwidth}
      \begin{definicao}[Arquitetura]
        As grandes estruturas de um sistema: seus elementos, as relações entre eles e as decisões \textbf{difíceis de mudar depois}. Ela existe para proteger atributos de qualidade.
      \end{definicao}
      \scriptsize
      \begin{tabularx}{\linewidth}{@{}l L@{}}
        \toprule
        \cabecalho{Atributo} & \cabecalho{Pergunta no LabMonitor} \\
        \midrule
        Manutenibilidade & trocar a API exige mexer nas telas? \\
        Testabilidade & dá para testar a regra de reserva sem um aparelho? \\
        Desempenho & a lista de leituras rola sem engasgar? \\
        Funcionamento \emph{offline} & o app mostra dados sem internet? \\
        Segurança & o \emph{token} fica protegido? \\
        \bottomrule
      \end{tabularx}
    \end{column}
    \begin{column}{0.46\textwidth}
      \begin{analogia}
        A planta de uma casa decide por onde passam os canos antes da primeira parede. Pintar um quarto não é arquitetura; mudar a cozinha de lugar é.
      \end{analogia}
      \begin{conceito}[Boa arquitetura em um app]
        Permite mudar a interface sem tocar nos dados, trocar a API sem tocar na interface e testar regras sem abrir o app.
      \end{conceito}
      \begin{dica}
        Coesão, acoplamento, SOLID e padrões de projeto foram vistos no Módulo 03: aqui eles ganham escala.
      \end{dica}
    \end{column}
  \end{columns}
\end{frame}

\begin{frame}{Padrões de apresentação: de MVC a MVI}
  \centering
  \begin{tikzpicture}[every node/.style={font=\tiny}]
    \draw[seta] (0,0) -- (13.4,0);
    \foreach \x/\t/\d in {1.2/MVC/{1979, Smalltalk\\ (Trygve Reenskaug)},4.6/MVP/{anos 1990 (Taligent);\\ comum no Android até 2017},8.1/MVVM/{2005, Microsoft (WPF);\\ padrão no Android e no Flutter},11.6/MVI/{anos 2010, inspirado\\ no Redux e no Cycle.js}}{
      \fill[ambar] (\x,0) circle (2pt);
      \node[font=\scriptsize\bfseries,text=azulpetroleo,above=1mm] at (\x,0) {\t};
      \node[rotulo,below=1mm] at (\x,0) {\d};
    }
  \end{tikzpicture}
  \raggedright
  \scriptsize
  \begin{tabularx}{\textwidth}{@{}l L L L@{}}
    \toprule
    \cabecalho{Padrão} & \cabecalho{Quem atualiza a tela} & \cabecalho{Ponto forte} & \cabecalho{Ponto fraco} \\
    \midrule
    MVC & o \emph{controller} manipula a \emph{view} & simples & no Android, a \cod{Activity} virava um \emph{controller} gigante \\
    MVP & o \emph{presenter} chama métodos de uma interface da \emph{view} & testável sem a tela & muito código de interface; guarda referência à \emph{view} \\
    MVVM & a \emph{view} observa o estado exposto pelo \emph{view model} & casa com UI declarativa e sobrevive à rotação & vários fluxos de estado podem se desencontrar \\
    MVI & estado único e imutável; cada intenção gera um novo estado & previsível e fácil de depurar & cerimônia demais para telas simples \\
    \bottomrule
  \end{tabularx}
  \normalsize
  \begin{conceito}[Na prática atual]
    Compose e Flutter usam MVVM com um estado único por tela (\cod{UiState}): uma mistura de MVVM com a disciplina do MVI.
  \end{conceito}
\end{frame}

\begin{frame}{Arquitetura limpa e arquitetura hexagonal}
  \begin{columns}[c]
    \begin{column}{0.46\textwidth}
      \centering
      \begin{tikzpicture}[every node/.style={font=\tiny,align=center}]
        \fill[cinzaclaro] (0,0) circle (2.15);
        \fill[cinzafundo,draw=azulpetroleo!40] (0,0) circle (1.6);
        \fill[ambarclaro!60,draw=ambar] (0,0) circle (1.05);
        \fill[azulpetroleo] (0,0) circle (0.55);
        \node[text=white,font=\tiny\bfseries] at (0,0) {Entidades};
        \node at (0,0.8) {Casos de uso};
        \node at (0,1.33) {Adaptadores};
        \node at (0,1.88) {Frameworks e drivers};
        \node[rotulo] at (0,-1.33) {\emph{view models}, repositórios};
        \node[rotulo] at (0,-1.88) {UI, banco, HTTP, BLE};
        \draw[setaambar] (2.6,-1.2) -- (1.2,-0.55);
        \node[rotulo,text width=1.5cm] at (3.05,-1.5) {dependências apontam para dentro};
      \end{tikzpicture}
    \end{column}
    \begin{column}{0.52\textwidth}
      \begin{definicao}[Regra da dependência]
        O código das camadas internas não conhece as externas: as regras de negócio não importam nada de Android, Flutter, HTTP ou SQL.
      \end{definicao}
      \begin{historia}
        Alistair Cockburn propôs a arquitetura hexagonal, de portas e adaptadores, em 2005. Robert C. Martin reuniu ideias parecidas na Arquitetura Limpa, em 2012.
      \end{historia}
      \begin{dica}[Na medida certa]
        Em apps pequenos, camadas de UI e de dados bastam. Crie a camada de domínio quando regras se repetirem ou precisarem de testes isolados, como recomendam o Android e o Flutter.
      \end{dica}
    \end{column}
  \end{columns}
\end{frame}

\begin{frame}[fragile]{Inversão de dependência na prática}
  \begin{columns}[T]
    \begin{column}{0.5\textwidth}
      \centering
      \begin{tikzpicture}[every node/.style={font=\tiny}]
        \node[rotulo,anchor=west] at (-1.1,0.55) {\textbf{sem inversão}};
        \node[caixaescura,font=\tiny,minimum width=2cm] (r1) at (0,0) {regras de leitura};
        \node[caixa,font=\tiny,minimum width=1.8cm] (b1) at (3.4,0) {\cod{FonteBle}};
        \draw[seta] (r1) -- node[rotulo,above]{conhece} (b1);
        \node[rotulo,anchor=west] at (-1.1,-0.75) {\textbf{com inversão}};
        \node[caixaescura,font=\tiny,minimum width=2cm] (r2) at (0,-1.3) {regras de leitura};
        \node[caixaambar,font=\tiny,minimum width=2.4cm] (i2) at (3.4,-1.3) {\emph{interface}\\ \cod{FonteLeituras}};
        \node[caixa,font=\tiny,minimum width=1.35cm] (b2) at (1.95,-2.75) {\cod{FonteBle}};
        \node[caixa,font=\tiny,minimum width=1.35cm] (m2) at (3.4,-2.75) {\cod{FonteMqtt}};
        \node[caixa,font=\tiny,minimum width=1.35cm] (f2) at (4.85,-2.75) {\cod{FonteFake}};
        \draw[seta] (r2) -- node[rotulo,above]{usa} (i2);
        \foreach \n in {b2,m2,f2}{\draw[setaambar] (\n.north) -- (\n.north |- i2.south);}
        \node[rotulo,below=1mm of m2] {implementam a interface};
      \end{tikzpicture}
\begin{codigo}[title=Dart]{Dart}
abstract interface class FonteLeituras {
  Stream<Leitura> leituras();
}
\end{codigo}
    \end{column}
    \begin{column}{0.48\textwidth}
      \begin{definicao}[Inversão de dependência]
        Módulos de alto nível não dependem de detalhes; ambos dependem de abstrações. Quem precisa define a interface; a camada de dados a implementa.
      \end{definicao}
\begin{codigo}[title=Kotlin]{Kotlin}
interface FonteLeituras {
  fun leituras(): Flow<Leitura>
}
\end{codigo}
      \begin{conceito}[O que se ganha]
        Trocar BLE por MQTT sem mexer nas regras, testar com uma fonte falsa e fazer a tela antes de o hardware existir.
      \end{conceito}
    \end{column}
  \end{columns}
\end{frame}

\begin{frame}{Injeção de dependência e localizador de serviços}
  \begin{columns}[T]
    \begin{column}{0.5\textwidth}
      \begin{definicao}[Inversão de controle]
        A classe não cria o que usa: alguém de fora monta os objetos e os entrega. É o princípio de Hollywood: ``não nos chame, nós chamamos você''.
      \end{definicao}
      \centering
      \begin{tikzpicture}[every node/.style={font=\tiny}]
        \node[caixaambarforte,font=\tiny\bfseries,minimum width=5.6cm] (raiz) at (0,0) {raiz de composição (\cod{Application}, \cod{main})};
        \node[caixa,font=\tiny,minimum width=1.7cm] (h) at (-1.9,-1) {cliente HTTP};
        \node[caixa,font=\tiny,minimum width=1.7cm] (b) at (0,-1) {banco};
        \node[caixapetroleo,font=\tiny,minimum width=1.9cm] (r) at (-0.95,-2) {repositório};
        \node[caixaescura,font=\tiny,minimum width=1.7cm] (v) at (1.9,-2) {\emph{view model}};
        \draw[setatrac] (raiz.south -| h) -- (h); \draw[setatrac] (raiz.south -| b) -- (b);
        \draw[seta] (h) -- (r); \draw[seta] (b) -- (r);
        \draw[seta] (r) -- node[rotulo,above]{injeta} (v);
        \node[rotulo] at (1.9,-1.05) {tracejado: cria};
      \end{tikzpicture}
      \raggedright
      \begin{historia}
        Martin Fowler batizou o padrão de injeção de dependência em 2004.
      \end{historia}
    \end{column}
    \begin{column}{0.48\textwidth}
      \scriptsize
      \begin{tabularx}{\linewidth}{@{}l L@{}}
        \toprule
        \cabecalho{Abordagem} & \cabecalho{Como funciona} \\
        \midrule
        Construtor & recebe tudo pronto; dependências explícitas e fáceis de trocar em testes \\
        Contêiner & uma biblioteca monta o grafo: Hilt, Koin; Provider, Riverpod \\
        Localizador & a classe pede a um registro global (\cod{get\_it}): prático, mas esconde dependências e só falha em execução \\
        \bottomrule
      \end{tabularx}
      \normalsize
      \begin{pordentro}[Tempo de vida (escopo)]
        \textbf{Único no app}: cliente HTTP, banco. \textbf{Por tela}: \emph{view model}. \textbf{Por uso}: objeto novo a cada pedido. Escopo errado causa vazamento ou estado compartilhado por engano.
      \end{pordentro}
    \end{column}
  \end{columns}
\end{frame}

\begin{frame}{Decisões de arquitetura: trocas e registros}
  \begin{columns}[T]
    \begin{column}{0.5\textwidth}
      \begin{conceito}[Toda decisão é uma troca]
        Mais camadas trazem isolamento e testabilidade, mas também mais arquivos e indireção. Pese o tamanho do time, a vida esperada do app e o custo de errar.
      \end{conceito}
      \scriptsize
      \begin{tabularx}{\linewidth}{@{}l L@{}}
        \toprule
        \cabecalho{Princípio} & \cabecalho{Significado} \\
        \midrule
        YAGNI & não construa o que ainda não é necessário \\
        KISS & prefira a solução mais simples que funciona \\
        DRY & cada conhecimento em um só lugar \\
        Regra de três & abstraia na terceira repetição, não na primeira \\
        \bottomrule
      \end{tabularx}
    \end{column}
    \begin{column}{0.48\textwidth}
      \begin{definicao}[ADR (\emph{Architecture Decision Record})]
        Texto curto, versionado com o código, que registra uma decisão: contexto, opções, decisão e consequências. Responde no futuro: ``por que fizemos assim?''.
      \end{definicao}
      \begin{exemplo}[ADR 007 --- Estado das telas]
        \textbf{Contexto}: leituras do sensor em tempo real. \textbf{Decisão}: um \cod{UiState} imutável por tela, exposto como fluxo observável. \textbf{Consequência}: telas testáveis sem interface.
      \end{exemplo}
      \begin{historia}
        Michael Nygard propôs os ADRs em 2011.
      \end{historia}
    \end{column}
  \end{columns}
\end{frame}

%  (coesão, acoplamento, SOLID e padrões de projeto estão no Módulo 03)
% =============================================================================
\section{Fundamentos: arquitetura de software}

\begin{frame}{O que é arquitetura de software}
  \begin{columns}[T]
    \begin{column}{0.52\textwidth}
      \begin{definicao}[Arquitetura]
        As grandes estruturas de um sistema: seus elementos, as relações entre eles e as decisões \textbf{difíceis de mudar depois}. Ela existe para proteger atributos de qualidade.
      \end{definicao}
      \scriptsize
      \begin{tabularx}{\linewidth}{@{}l L@{}}
        \toprule
        \cabecalho{Atributo} & \cabecalho{Pergunta no LabMonitor} \\
        \midrule
        Manutenibilidade & trocar a API exige mexer nas telas? \\
        Testabilidade & dá para testar a regra de reserva sem um aparelho? \\
        Desempenho & a lista de leituras rola sem engasgar? \\
        Funcionamento \emph{offline} & o app mostra dados sem internet? \\
        Segurança & o \emph{token} fica protegido? \\
        \bottomrule
      \end{tabularx}
    \end{column}
    \begin{column}{0.46\textwidth}
      \begin{analogia}
        A planta de uma casa decide por onde passam os canos antes da primeira parede. Pintar um quarto não é arquitetura; mudar a cozinha de lugar é.
      \end{analogia}
      \begin{conceito}[Boa arquitetura em um app]
        Permite mudar a interface sem tocar nos dados, trocar a API sem tocar na interface e testar regras sem abrir o app.
      \end{conceito}
      \begin{dica}
        Coesão, acoplamento, SOLID e padrões de projeto foram vistos no Módulo 03: aqui eles ganham escala.
      \end{dica}
    \end{column}
  \end{columns}
\end{frame}

\begin{frame}{Padrões de apresentação: de MVC a MVI}
  \centering
  \begin{tikzpicture}[every node/.style={font=\tiny}]
    \draw[seta] (0,0) -- (13.4,0);
    \foreach \x/\t/\d in {1.2/MVC/{1979, Smalltalk\\ (Trygve Reenskaug)},4.6/MVP/{anos 1990 (Taligent);\\ comum no Android até 2017},8.1/MVVM/{2005, Microsoft (WPF);\\ padrão no Android e no Flutter},11.6/MVI/{anos 2010, inspirado\\ no Redux e no Cycle.js}}{
      \fill[ambar] (\x,0) circle (2pt);
      \node[font=\scriptsize\bfseries,text=azulpetroleo,above=1mm] at (\x,0) {\t};
      \node[rotulo,below=1mm] at (\x,0) {\d};
    }
  \end{tikzpicture}
  \raggedright
  \scriptsize
  \begin{tabularx}{\textwidth}{@{}l L L L@{}}
    \toprule
    \cabecalho{Padrão} & \cabecalho{Quem atualiza a tela} & \cabecalho{Ponto forte} & \cabecalho{Ponto fraco} \\
    \midrule
    MVC & o \emph{controller} manipula a \emph{view} & simples & no Android, a \cod{Activity} virava um \emph{controller} gigante \\
    MVP & o \emph{presenter} chama métodos de uma interface da \emph{view} & testável sem a tela & muito código de interface; guarda referência à \emph{view} \\
    MVVM & a \emph{view} observa o estado exposto pelo \emph{view model} & casa com UI declarativa e sobrevive à rotação & vários fluxos de estado podem se desencontrar \\
    MVI & estado único e imutável; cada intenção gera um novo estado & previsível e fácil de depurar & cerimônia demais para telas simples \\
    \bottomrule
  \end{tabularx}
  \normalsize
  \begin{conceito}[Na prática atual]
    Compose e Flutter usam MVVM com um estado único por tela (\cod{UiState}): uma mistura de MVVM com a disciplina do MVI.
  \end{conceito}
\end{frame}

\begin{frame}{Arquitetura limpa e arquitetura hexagonal}
  \begin{columns}[c]
    \begin{column}{0.46\textwidth}
      \centering
      \begin{tikzpicture}[every node/.style={font=\tiny,align=center}]
        \fill[cinzaclaro] (0,0) circle (2.15);
        \fill[cinzafundo,draw=azulpetroleo!40] (0,0) circle (1.6);
        \fill[ambarclaro!60,draw=ambar] (0,0) circle (1.05);
        \fill[azulpetroleo] (0,0) circle (0.55);
        \node[text=white,font=\tiny\bfseries] at (0,0) {Entidades};
        \node at (0,0.8) {Casos de uso};
        \node at (0,1.33) {Adaptadores};
        \node at (0,1.88) {Frameworks e drivers};
        \node[rotulo] at (0,-1.33) {\emph{view models}, repositórios};
        \node[rotulo] at (0,-1.88) {UI, banco, HTTP, BLE};
        \draw[setaambar] (2.6,-1.2) -- (1.2,-0.55);
        \node[rotulo,text width=1.5cm] at (3.05,-1.5) {dependências apontam para dentro};
      \end{tikzpicture}
    \end{column}
    \begin{column}{0.52\textwidth}
      \begin{definicao}[Regra da dependência]
        O código das camadas internas não conhece as externas: as regras de negócio não importam nada de Android, Flutter, HTTP ou SQL.
      \end{definicao}
      \begin{historia}
        Alistair Cockburn propôs a arquitetura hexagonal, de portas e adaptadores, em 2005. Robert C. Martin reuniu ideias parecidas na Arquitetura Limpa, em 2012.
      \end{historia}
      \begin{dica}[Na medida certa]
        Em apps pequenos, camadas de UI e de dados bastam. Crie a camada de domínio quando regras se repetirem ou precisarem de testes isolados, como recomendam o Android e o Flutter.
      \end{dica}
    \end{column}
  \end{columns}
\end{frame}

\begin{frame}[fragile]{Inversão de dependência na prática}
  \begin{columns}[T]
    \begin{column}{0.5\textwidth}
      \centering
      \begin{tikzpicture}[every node/.style={font=\tiny}]
        \node[rotulo,anchor=west] at (-1.1,0.55) {\textbf{sem inversão}};
        \node[caixaescura,font=\tiny,minimum width=2cm] (r1) at (0,0) {regras de leitura};
        \node[caixa,font=\tiny,minimum width=1.8cm] (b1) at (3.4,0) {\cod{FonteBle}};
        \draw[seta] (r1) -- node[rotulo,above]{conhece} (b1);
        \node[rotulo,anchor=west] at (-1.1,-0.75) {\textbf{com inversão}};
        \node[caixaescura,font=\tiny,minimum width=2cm] (r2) at (0,-1.3) {regras de leitura};
        \node[caixaambar,font=\tiny,minimum width=2.4cm] (i2) at (3.4,-1.3) {\emph{interface}\\ \cod{FonteLeituras}};
        \node[caixa,font=\tiny,minimum width=1.35cm] (b2) at (1.95,-2.75) {\cod{FonteBle}};
        \node[caixa,font=\tiny,minimum width=1.35cm] (m2) at (3.4,-2.75) {\cod{FonteMqtt}};
        \node[caixa,font=\tiny,minimum width=1.35cm] (f2) at (4.85,-2.75) {\cod{FonteFake}};
        \draw[seta] (r2) -- node[rotulo,above]{usa} (i2);
        \foreach \n in {b2,m2,f2}{\draw[setaambar] (\n.north) -- (\n.north |- i2.south);}
        \node[rotulo,below=1mm of m2] {implementam a interface};
      \end{tikzpicture}
\begin{codigo}[title=Dart]{Dart}
abstract interface class FonteLeituras {
  Stream<Leitura> leituras();
}
\end{codigo}
    \end{column}
    \begin{column}{0.48\textwidth}
      \begin{definicao}[Inversão de dependência]
        Módulos de alto nível não dependem de detalhes; ambos dependem de abstrações. Quem precisa define a interface; a camada de dados a implementa.
      \end{definicao}
\begin{codigo}[title=Kotlin]{Kotlin}
interface FonteLeituras {
  fun leituras(): Flow<Leitura>
}
\end{codigo}
      \begin{conceito}[O que se ganha]
        Trocar BLE por MQTT sem mexer nas regras, testar com uma fonte falsa e fazer a tela antes de o hardware existir.
      \end{conceito}
    \end{column}
  \end{columns}
\end{frame}

\begin{frame}{Injeção de dependência e localizador de serviços}
  \begin{columns}[T]
    \begin{column}{0.5\textwidth}
      \begin{definicao}[Inversão de controle]
        A classe não cria o que usa: alguém de fora monta os objetos e os entrega. É o princípio de Hollywood: ``não nos chame, nós chamamos você''.
      \end{definicao}
      \centering
      \begin{tikzpicture}[every node/.style={font=\tiny}]
        \node[caixaambarforte,font=\tiny\bfseries,minimum width=5.6cm] (raiz) at (0,0) {raiz de composição (\cod{Application}, \cod{main})};
        \node[caixa,font=\tiny,minimum width=1.7cm] (h) at (-1.9,-1) {cliente HTTP};
        \node[caixa,font=\tiny,minimum width=1.7cm] (b) at (0,-1) {banco};
        \node[caixapetroleo,font=\tiny,minimum width=1.9cm] (r) at (-0.95,-2) {repositório};
        \node[caixaescura,font=\tiny,minimum width=1.7cm] (v) at (1.9,-2) {\emph{view model}};
        \draw[setatrac] (raiz.south -| h) -- (h); \draw[setatrac] (raiz.south -| b) -- (b);
        \draw[seta] (h) -- (r); \draw[seta] (b) -- (r);
        \draw[seta] (r) -- node[rotulo,above]{injeta} (v);
        \node[rotulo] at (1.9,-1.05) {tracejado: cria};
      \end{tikzpicture}
      \raggedright
      \begin{historia}
        Martin Fowler batizou o padrão de injeção de dependência em 2004.
      \end{historia}
    \end{column}
    \begin{column}{0.48\textwidth}
      \scriptsize
      \begin{tabularx}{\linewidth}{@{}l L@{}}
        \toprule
        \cabecalho{Abordagem} & \cabecalho{Como funciona} \\
        \midrule
        Construtor & recebe tudo pronto; dependências explícitas e fáceis de trocar em testes \\
        Contêiner & uma biblioteca monta o grafo: Hilt, Koin; Provider, Riverpod \\
        Localizador & a classe pede a um registro global (\cod{get\_it}): prático, mas esconde dependências e só falha em execução \\
        \bottomrule
      \end{tabularx}
      \normalsize
      \begin{pordentro}[Tempo de vida (escopo)]
        \textbf{Único no app}: cliente HTTP, banco. \textbf{Por tela}: \emph{view model}. \textbf{Por uso}: objeto novo a cada pedido. Escopo errado causa vazamento ou estado compartilhado por engano.
      \end{pordentro}
    \end{column}
  \end{columns}
\end{frame}

\begin{frame}{Decisões de arquitetura: trocas e registros}
  \begin{columns}[T]
    \begin{column}{0.5\textwidth}
      \begin{conceito}[Toda decisão é uma troca]
        Mais camadas trazem isolamento e testabilidade, mas também mais arquivos e indireção. Pese o tamanho do time, a vida esperada do app e o custo de errar.
      \end{conceito}
      \scriptsize
      \begin{tabularx}{\linewidth}{@{}l L@{}}
        \toprule
        \cabecalho{Princípio} & \cabecalho{Significado} \\
        \midrule
        YAGNI & não construa o que ainda não é necessário \\
        KISS & prefira a solução mais simples que funciona \\
        DRY & cada conhecimento em um só lugar \\
        Regra de três & abstraia na terceira repetição, não na primeira \\
        \bottomrule
      \end{tabularx}
    \end{column}
    \begin{column}{0.48\textwidth}
      \begin{definicao}[ADR (\emph{Architecture Decision Record})]
        Texto curto, versionado com o código, que registra uma decisão: contexto, opções, decisão e consequências. Responde no futuro: ``por que fizemos assim?''.
      \end{definicao}
      \begin{exemplo}[ADR 007 --- Estado das telas]
        \textbf{Contexto}: leituras do sensor em tempo real. \textbf{Decisão}: um \cod{UiState} imutável por tela, exposto como fluxo observável. \textbf{Consequência}: telas testáveis sem interface.
      \end{exemplo}
      \begin{historia}
        Michael Nygard propôs os ADRs em 2011.
      \end{historia}
    \end{column}
  \end{columns}
\end{frame}


% =============================================================================
\section{Por que arquitetura?}

\begin{frame}{O problema da ``Activity faz-tudo''}
  \begin{columns}[c]
    \begin{column}{0.46\textwidth}
      \centering
      \begin{tikzpicture}
        \node[caixaescura,minimum width=4.4cm,minimum height=3.6cm,align=center] (a) {};
        \node[font=\small\bfseries,text=white] at ([yshift=-3mm]a.north) {MainActivity (2\,000 linhas)};
        \foreach \t/\y in {UI e layout/0.85,chamadas HTTP/0.4,SQL e arquivos/-0.05,regras de negócio/-0.5,Bluetooth/-0.95}{
          \node[caixaambar,minimum width=3.4cm,minimum height=4mm,font=\tiny] at (0,\y) {\t};}
      \end{tikzpicture}
    \end{column}
    \begin{column}{0.52\textwidth}
      \begin{erro}[Consequências]
        \begin{itemize}
          \item Rotação de tela refaz requisições e perde dados.
          \item Impossível testar regras sem abrir o app.
          \item Uma mudança na API quebra a tela inteira.
        \end{itemize}
      \end{erro}
      \begin{conceito}[Princípios recomendados pelo Google]
        \begin{itemize}
          \item Separação de responsabilidades.
          \item UI guiada por modelos de dados.
          \item Fonte única da verdade.
          \item Fluxo unidirecional de dados (UDF).
        \end{itemize}
      \end{conceito}
    \end{column}
  \end{columns}
\end{frame}

\begin{frame}{A arquitetura em camadas}
  \begin{columns}[c]
    \begin{column}{0.58\textwidth}
      \centering
      \begin{tikzpicture}[node distance=3mm]
        \node[caixa,minimum width=2.9cm,minimum height=8mm] (ui) {Elementos de UI\\ \tiny composables};
        \node[caixaescura,minimum width=2.9cm,minimum height=8mm,right=4mm of ui] (vm) {State holders\\ \tiny ViewModel};
        \node[pontilhado,fit=(ui)(vm),label={[rotulo,anchor=south west]north west:\textbf{Camada de UI}}] (cui) {};
        \node[caixaambar,minimum width=6.2cm,minimum height=7mm,below=8mm of cui] (uc) {Casos de uso \tiny (opcional)};
        \node[pontilhado,fit=(uc),label={[rotulo,anchor=south west]north west:\textbf{Camada de domínio}}] (cdom) {};
        \node[caixapetroleo,minimum width=6.2cm,minimum height=7mm,below=8mm of cdom] (repo) {Repositórios};
        \node[caixa,minimum width=2.9cm,below=3mm of repo.south west,anchor=north west] (loc) {Fonte local\\ \tiny Room, DataStore};
        \node[caixa,minimum width=2.9cm,below=3mm of repo.south east,anchor=north east] (rem) {Fonte remota\\ \tiny API, Firebase, BLE};
        \node[pontilhado,fit=(repo)(loc)(rem),label={[rotulo,anchor=south west]north west:\textbf{Camada de dados}}] (cdad) {};
        \draw[seta] (ui) -- (vm);
        \draw[seta] (vm.south) -- ++(0,-4mm) -| ([xshift=10mm]uc.north);
        \draw[seta] (uc) -- (repo);
        \draw[seta] (repo.south -| loc) -- (loc);
        \draw[seta] (repo.south -| rem) -- (rem);
      \end{tikzpicture}
    \end{column}
    \begin{column}{0.4\textwidth}
      \small
      \begin{itemize}
        \item As setas indicam \textbf{dependência}: cada camada só conhece a de baixo.
        \item A camada de dados \textbf{não sabe} que existe UI.
        \item Os dados sobem como \cod{Flow}; os eventos descem como chamadas de função.
      \end{itemize}
      \begin{dica}
        Comece com UI + dados. Crie a camada de domínio quando regras se repetirem em mais de um ViewModel.
      \end{dica}
    \end{column}
  \end{columns}
\end{frame}

% =============================================================================
\section{Camada de UI: MVVM}

\begin{frame}{MVVM e fluxo unidirecional}
  \begin{center}
  \begin{tikzpicture}[node distance=24mm]
    \node[caixa,minimum width=2.8cm,minimum height=11mm] (v) {\textbf{View}\\ \tiny composables (stateless)};
    \node[caixaescura,minimum width=2.8cm,minimum height=11mm,right=of v] (vm) {\textbf{ViewModel}\\ \tiny estado + lógica de tela};
    \node[caixapetroleo,minimum width=2.8cm,minimum height=11mm,right=of vm] (m) {\textbf{Model}\\ \tiny repositórios e dados};
    \draw[setaambar] ([yshift=2mm]v.east) -- node[rotulo,above]{eventos: \cod{onReservar()}} ([yshift=2mm]vm.west);
    \draw[seta] ([yshift=-2mm]vm.west) -- node[rotulo,below]{\cod{StateFlow<UiState>}} ([yshift=-2mm]v.east);
    \draw[setaambar] ([yshift=2mm]vm.east) -- node[rotulo,above]{chamadas} ([yshift=2mm]m.west);
    \draw[seta] ([yshift=-2mm]m.west) -- node[rotulo,below]{\cod{Flow<dados>}} ([yshift=-2mm]vm.east);
  \end{tikzpicture}
  \end{center}
  \begin{columns}[T]
    \begin{column}{0.32\textwidth}
      \begin{conceito}[View]
        Desenha o \cod{UiState} e repassa eventos. Não decide nada.
      \end{conceito}
    \end{column}
    \begin{column}{0.32\textwidth}
      \begin{conceito}[ViewModel]
        Sobrevive à rotação, transforma dados em \cod{UiState} e trata eventos.
      \end{conceito}
    \end{column}
    \begin{column}{0.32\textwidth}
      \begin{conceito}[Model]
        Dados e regras do app, sem nenhuma dependência de Android UI.
      \end{conceito}
    \end{column}
  \end{columns}
  \begin{atencao}
    O ViewModel \textbf{nunca} guarda referência a \cod{Activity}, \cod{Context} de tela ou composables: ele vive mais que eles e causaria vazamento de memória.
  \end{atencao}
\end{frame}

\begin{frame}[fragile]{Modelando o \cod{UiState}}
  \begin{columns}[T]
    \begin{column}{0.5\textwidth}
\begin{codigo}[title=Estado como data class]{Kotlin}
data class PainelUiState(
  val carregando: Boolean = true,
  val bancadas: List<BancadaUi> = emptyList(),
  val filtro: String = "",
  val mensagem: String? = null    // ex.: Snackbar
) {
  val vazio get() = !carregando && bancadas.isEmpty()
}

data class BancadaUi(
  val id: Int,
  val titulo: String,
  val temperatura: String,       // já formatada
  val emAlerta: Boolean
)
\end{codigo}
    \end{column}
    \begin{column}{0.48\textwidth}
\begin{codigo}[title=Ou como sealed interface]{Kotlin}
sealed interface DetalheUiState {
  data object Carregando : DetalheUiState
  data class Pronto(val eq: Equipamento)
    : DetalheUiState
  data class Erro(val msg: String)
    : DetalheUiState
}
\end{codigo}
      \begin{dica}[Qual escolher?]
        \textbf{sealed}: estados exclusivos (carregando, pronto ou erro). \textbf{data class}: aspectos que convivem (lista, filtro, mensagem).
      \end{dica}
      \begin{conceito}
        O \cod{UiState} é \textbf{imutável}: cada mudança gera uma cópia com \cod{copy()}.
      \end{conceito}
    \end{column}
  \end{columns}
\end{frame}

\begin{frame}[fragile]{O ViewModel}
  \begin{columns}[T]
    \begin{column}{0.58\textwidth}
\begin{codigo}[title=PainelViewModel.kt]{Kotlin}
class PainelViewModel(
  private val repo: BancadaRepository
) : ViewModel() {
  private val filtro = MutableStateFlow("")
  val uiState: StateFlow<PainelUiState> =
    combine(repo.bancadas(), filtro) { lista, f ->
      PainelUiState(
        carregando = false,
        bancadas = lista
          .filter { f.isBlank() || f in it.nome }
          .map { it.paraUi() },
        filtro = f
      )
    }.stateIn(viewModelScope,
      SharingStarted.WhileSubscribed(5_000),
      PainelUiState())

  fun onFiltroChange(novo: String) { filtro.value = novo }
  fun onAtualizar() =
    viewModelScope.launch { repo.sincronizar() }
}
\end{codigo}
    \end{column}
    \begin{column}{0.4\textwidth}
      \small
      \begin{itemize}
        \item Entradas: fluxos do repositório e estado local (\cod{filtro}).
        \item \cod{combine} gera um novo \cod{UiState} sempre que algo muda.
        \item \cod{stateIn} entrega um \cod{StateFlow} pronto para a UI.
        \item Eventos são funções públicas com nomes de intenção (\cod{onAtualizar}).
      \end{itemize}
      \begin{dica}
        \cod{paraUi()} é uma função de mapeamento: formata números e datas \textbf{fora} da UI, o que facilita testar.
      \end{dica}
    \end{column}
  \end{columns}
\end{frame}

\begin{frame}[fragile]{Tela: \emph{route} + conteúdo sem estado}
  \begin{columns}[T]
    \begin{column}{0.54\textwidth}
\begin{codigo}[title=TelaPainel.kt]{Kotlin}
@Composable  // liga o ViewModel à tela (sem lógica)
fun PainelRoute(
  onAbrir: (Int) -> Unit,
  vm: PainelViewModel = hiltViewModel()
) {
  val estado by vm.uiState.collectAsStateWithLifecycle()
  TelaPainel(
    estado = estado,
    onFiltroChange = vm::onFiltroChange,
    onAtualizar = vm::onAtualizar,
    onAbrir = onAbrir
  )
}
@Composable  // só desenha: fácil de testar
fun TelaPainel(
  estado: PainelUiState,
  onFiltroChange: (String) -> Unit,
  onAtualizar: () -> Unit,
  onAbrir: (Int) -> Unit
) { /* Scaffold, campo de busca, LazyColumn... */ }
\end{codigo}
    \end{column}
    \begin{column}{0.44\textwidth}
      \begin{conceito}[Separação Route / Screen]
        \textbf{Route}: obtém o ViewModel e coleta o estado.\\
        \textbf{Screen}: recebe estado e lambdas; não conhece o ViewModel.
      \end{conceito}
      \begin{dica}[Preview com qualquer estado]
        \cod{TelaPainel(PainelUiState(carregando = false,}\\
        \cod{\ \ bancadas = exemplo), \{\}, \{\}, \{\})}\\[2pt]
        Pré-visualize vazio, carregando, erro e lista cheia sem rodar o app.
      \end{dica}
    \end{column}
  \end{columns}
\end{frame}

% =============================================================================
\section{Camada de dados}

\begin{frame}[fragile]{Repositório: fonte única da verdade}
  \begin{columns}[T]
    \begin{column}{0.52\textwidth}
\begin{codigo}[title=BancadaRepository.kt]{Kotlin}
interface BancadaRepository {
  fun bancadas(): Flow<List<Bancada>>
  suspend fun sincronizar()
  suspend fun reservar(id: Int, ate: Instant)
}

class BancadaRepositoryImpl(
  private val dao: BancadaDao,        // local (M10)
  private val api: LabApi             // remota (M11)
) : BancadaRepository {
  // a UI sempre lê do banco local
  override fun bancadas() = dao.observarTodas()
    .map { lista -> lista.map { it.paraModelo() } }
  // a rede só atualiza o banco
  override suspend fun sincronizar() {
    val remotas = api.listarBancadas()
    dao.substituirTodas(remotas.map { it.paraEntidade() })
  }
  override suspend fun reservar(id: Int, ate: Instant) =
    api.reservar(id, ate)
}
\end{codigo}
    \end{column}
    \begin{column}{0.46\textwidth}
      \centering
      \begin{tikzpicture}[node distance=6mm]
        \node[caixaescura,minimum width=3.6cm] (vm) {ViewModel};
        \node[caixapetroleo,minimum width=3.6cm,below=of vm] (r) {Repositório};
        \node[caixaambar,below left=6mm and -14mm of r] (db) {Room\\ \tiny fonte da verdade};
        \node[caixa,below right=6mm and -14mm of r] (api) {API\\ \tiny rede};
        \draw[seta] (db) -- node[rotulo,left]{Flow} (r);
        \draw[seta] (r) -- (vm);
        \draw[setaambar] (api) -- node[rotulo,below,sloped]{grava} (db);
      \end{tikzpicture}
      \begin{conceito}[Por que interface?]
        A UI depende da \textbf{abstração}. Nos testes trocamos a implementação por um \emph{fake} em memória.
      \end{conceito}
    \end{column}
  \end{columns}
\end{frame}

\begin{frame}[fragile]{Um modelo para cada camada}
  \begin{columns}[T]
    \begin{column}{0.5\textwidth}
\begin{codigo}[title=Modelos e mapeamentos]{Kotlin}
// Rede (DTO): espelha o JSON da API
@Serializable
data class BancadaDto(
  val id: Int, val nome: String,
  @SerialName("temp_c") val temp: Double?)

// Banco (Entity): espelha a tabela
@Entity data class BancadaEntity(
  @PrimaryKey val id: Int, val nome: String,
  val temp: Double?)

// Domínio: o que o app entende
data class Bancada(val id: Int, val nome: String,
  val temperatura: Double?)

fun BancadaDto.paraEntidade() =
  BancadaEntity(id, nome, temp)
fun BancadaEntity.paraModelo() =
  Bancada(id, nome, temp)
\end{codigo}
    \end{column}
    \begin{column}{0.48\textwidth}
      \centering
      \begin{tikzpicture}[node distance=5mm]
        \node[caixa] (dto) {\cod{BancadaDto}};
        \node[caixaambar,right=of dto] (ent) {\cod{BancadaEntity}};
        \node[caixapetroleo,below=of ent] (mod) {\cod{Bancada}};
        \node[caixaescura,left=of mod] (ui) {\cod{BancadaUi}};
        \draw[seta] (dto) -- (ent); \draw[seta] (ent) -- (mod); \draw[seta] (mod) -- (ui);
      \end{tikzpicture}
      \begin{conceito}[Por que tantos modelos?]
        Se a API renomear um campo, só o DTO e um mapeamento mudam. O resto do app nem fica sabendo.
      \end{conceito}
      \begin{dica}
        Em apps pequenos é aceitável usar o mesmo modelo em duas camadas. Separe quando os formatos começarem a divergir.
      \end{dica}
    \end{column}
  \end{columns}
\end{frame}

% =============================================================================
\section{Camada de domínio}

\begin{frame}[fragile]{Casos de uso}
  \begin{columns}[T]
    \begin{column}{0.54\textwidth}
\begin{codigo}[title=VerificarAlertasUseCase.kt]{Kotlin}
class VerificarAlertasUseCase(
  private val bancadas: BancadaRepository,
  private val ajustes: AjustesRepository
) {
  // operator invoke: o objeto é chamado como função
  operator fun invoke(): Flow<List<Alerta>> =
    combine(bancadas.bancadas(), ajustes.limite()) {
      lista, limite ->
      lista.mapNotNull { b ->
        val t = b.temperatura ?: return@mapNotNull null
        if (t > limite) Alerta(b.id, t, limite) else null
      }
    }
}

// no ViewModel
val alertas = verificarAlertas()
  .stateIn(viewModelScope, WhileSubscribed(5_000),
    emptyList())
\end{codigo}
    \end{column}
    \begin{column}{0.44\textwidth}
      \begin{conceito}[Caso de uso]
        Classe com \textbf{uma única responsabilidade} de negócio, nomeada com verbo: \cod{ReservarBancada}, \cod{CalcularConsumo}, \cod{VerificarAlertas}.
      \end{conceito}
      \begin{dica}[Quando criar]
        Quando a regra é usada por vários ViewModels, combina repositórios ou o ViewModel cresceu demais.
      \end{dica}
      \begin{atencao}
        Domínio é Kotlin puro, sem \cod{Context} nem classes Android, e por isso é reaproveitável em Kotlin Multiplatform.
      \end{atencao}
    \end{column}
  \end{columns}
\end{frame}

% =============================================================================
\section{Injeção de dependência com Hilt}

\begin{frame}[fragile]{O que é injeção de dependência}
  \begin{columns}[T]
    \begin{column}{0.48\textwidth}
\begin{codigo}[title=Sem DI: a classe cria o que usa]{Kotlin}
class PainelViewModel : ViewModel() {
  private val repo = BancadaRepositoryImpl(
    AppDatabase.get(???).bancadaDao(),
    Retrofit.Builder()...create(LabApi::class)
  )  // acoplado e impossível de testar
}
\end{codigo}
\begin{codigo}[title=Com DI: recebe pronto]{Kotlin}
class PainelViewModel(
  private val repo: BancadaRepository
) : ViewModel()
\end{codigo}
    \end{column}
    \begin{column}{0.5\textwidth}
      \begin{conceito}[Injeção de dependência (DI)]
        Uma classe \textbf{recebe} suas dependências (pelo construtor) em vez de criá-las. Alguém de fora monta o ``grafo'' de objetos.
      \end{conceito}
      \begin{itemize}
        \item \textbf{Teste}: injete um repositório \emph{fake}.
        \item \textbf{Reuso}: um único banco e um único cliente HTTP para o app todo.
        \item \textbf{Troca}: mude a implementação em um só lugar.
      \end{itemize}
      \begin{dica}[Opções]
        \textbf{Hilt} (oficial Google, verificado em compilação), \textbf{Koin} (simples, também multiplataforma) ou DI manual com uma classe \cod{AppContainer}.
      \end{dica}
    \end{column}
  \end{columns}
\end{frame}

\begin{frame}[fragile]{Hilt: configuração}
  \begin{columns}[T]
    \begin{column}{0.5\textwidth}
\begin{codigo}[title=app/build.gradle.kts]{Kotlin}
plugins {
  alias(libs.plugins.ksp)
  alias(libs.plugins.hilt)
}
dependencies {
  implementation(libs.hilt.android)
  ksp(libs.hilt.compiler)
  implementation(libs.hilt.navigation.compose)
}
\end{codigo}
\begin{codigo}[title=LabMonitorApp.kt]{Kotlin}
@HiltAndroidApp
class LabMonitorApp : Application()
\end{codigo}
\begin{codigo}[title=MainActivity.kt]{Kotlin}
@AndroidEntryPoint
class MainActivity : ComponentActivity() { /*...*/ }
\end{codigo}
    \end{column}
    \begin{column}{0.48\textwidth}
      \begin{enumerate}
        \item Adicione os plugins e as dependências.
        \item Crie uma classe \cod{Application} anotada com \cod{@HiltAndroidApp} e registre-a no manifesto (\cod{android:name=".LabMonitorApp"}).
        \item Anote a Activity com \cod{@AndroidEntryPoint}.
        \item Anote ViewModels com \cod{@HiltViewModel}.
        \item Ensine o Hilt a criar interfaces e classes de bibliotecas com \cod{@Module}.
      \end{enumerate}
      \begin{dica}
        O \textbf{KSP} (\emph{Kotlin Symbol Processing}) gera o código de DI em tempo de compilação: erros de configuração aparecem no build, não no celular.
      \end{dica}
    \end{column}
  \end{columns}
\end{frame}

\begin{frame}[fragile]{Hilt: ViewModels e módulos}
  \begin{columns}[T]
    \begin{column}{0.5\textwidth}
\begin{codigo}[title=ViewModel injetado]{Kotlin}
@HiltViewModel
class PainelViewModel @Inject constructor(
  private val repo: BancadaRepository,
  private val verificarAlertas: VerificarAlertasUseCase
) : ViewModel() { /* ... */ }

// Classes próprias: basta @Inject no construtor
class VerificarAlertasUseCase @Inject constructor(
  private val bancadas: BancadaRepository,
  private val ajustes: AjustesRepository
)

// Na tela
val vm: PainelViewModel = hiltViewModel()
\end{codigo}
    \end{column}
    \begin{column}{0.48\textwidth}
\begin{codigo}[title=di/DataModule.kt]{Kotlin}
@Module
@InstallIn(SingletonComponent::class)
abstract class RepositoryModule {
  @Binds @Singleton   // interface -> implementação
  abstract fun bancadas(
    impl: BancadaRepositoryImpl
  ): BancadaRepository
}

@Module
@InstallIn(SingletonComponent::class)
object DatabaseModule {
  @Provides @Singleton  // classes de biblioteca
  fun banco(@ApplicationContext ctx: Context) =
    Room.databaseBuilder(ctx, AppDatabase::class.java,
      "labmonitor.db").build()

  @Provides fun dao(db: AppDatabase) = db.bancadaDao()
}
\end{codigo}
    \end{column}
  \end{columns}
\end{frame}

% =============================================================================
\section{Organização do projeto}

\begin{frame}[fragile]{Pacotes por funcionalidade e modularização}
  \begin{columns}[T]
    \begin{column}{0.44\textwidth}
\begin{codigo}[title=Pacotes por feature]{bash}
br.edu.ifma.labmonitor/
|-- LabMonitorApp.kt
|-- MainActivity.kt
|-- navegacao/
|-- core/
|   |-- data/   (db, rede, repositórios)
|   |-- domain/ (modelos, casos de uso)
|   `-- ui/     (tema, componentes comuns)
|-- feature/
|   |-- painel/
|   |   |-- PainelViewModel.kt
|   |   `-- TelaPainel.kt
|   |-- equipamentos/
|   `-- reservas/
`-- di/
\end{codigo}
    \end{column}
    \begin{column}{0.54\textwidth}
      \begin{conceito}[Pacotes por funcionalidade]
        Agrupar por \emph{feature} em vez de por tipo mantém juntos os arquivos que mudam juntos.
      \end{conceito}
      \centering
      \begin{tikzpicture}[node distance=2.5mm and 3mm]
        \node[caixaescura] (app) {\cod{:app}};
        \node[caixa,below left=of app] (f1) {\cod{:feature:painel}};
        \node[caixa,below right=of app] (f2) {\cod{:feature:reservas}};
        \node[caixaambar,below=of f1] (d) {\cod{:core:data}};
        \node[caixaambar,below=of f2] (u) {\cod{:core:ui}};
        \draw[seta] (app) -- (f1); \draw[seta] (app) -- (f2);
        \draw[seta] (f1) -- (d); \draw[seta] (f2) -- (d);
        \draw[seta] (f1) -- (u); \draw[seta] (f2) -- (u);
      \end{tikzpicture}
      \begin{dica}[Modularização Gradle]
        Em projetos grandes, cada pasta vira um módulo Gradle, com builds mais rápidos e fronteiras garantidas pelo compilador (veja o \emph{Now in Android}).
      \end{dica}
    \end{column}
  \end{columns}
\end{frame}

% =============================================================================
\section{Fechamento}

\begin{frame}{Resumo}
  \begin{resumo}
    \begin{itemize}
      \item Separe o app em camadas: UI (composables + ViewModel), domínio (casos de uso) e dados (repositórios + fontes).
      \item MVVM com UDF: o ViewModel expõe \cod{StateFlow<UiState>} e recebe eventos como funções.
      \item Divida cada tela em \emph{Route} (liga o ViewModel) e \emph{Screen} (só desenha).
      \item O repositório é a fonte única da verdade; a rede alimenta o banco local.
      \item DTO, Entity e modelo de domínio isolam o app de mudanças externas.
      \item Hilt monta o grafo de dependências: \cod{@HiltAndroidApp}, \cod{@HiltViewModel}, \cod{@Inject}, \cod{@Module}.
    \end{itemize}
  \end{resumo}
\end{frame}

\begin{frame}{Glossário (1/2)}
  \begin{glossario}
    Arquitetura & estruturas do sistema e decisões difíceis de mudar depois \\
    Camada & grupo de código com uma responsabilidade: UI, domínio ou dados \\
    MVVM & a \emph{view} observa o estado exposto pelo \cod{ViewModel} \\
    MVI & estado único e imutável; cada intenção do usuário gera um novo estado \\
    \cod{UiState} & classe imutável com tudo o que a tela precisa desenhar \\
    \cod{ViewModel} & guarda o estado da tela e sobrevive à rotação \\
    Repositório & fonte única da verdade de um tipo de dado \\
    Caso de uso & classe com uma única regra de negócio reutilizável \\
    Regra da dependência & camadas internas não conhecem as externas \\
  \end{glossario}
\end{frame}

\begin{frame}{Glossário (2/2)}
  \begin{glossario}
    Inversão de dependência & depender de abstrações, não de implementações \\
    Injeção de dependência & receber as dependências prontas em vez de criá-las \\
    Raiz de composição & único lugar onde os objetos concretos são criados e ligados \\
    Localizador de serviços & registro global ao qual a classe pede suas dependências \\
    Hilt & biblioteca de DI do Google, com o grafo verificado na compilação \\
    Escopo & tempo de vida de uma dependência (\cod{@Singleton}, \cod{@ViewModelScoped}) \\
    ADR & registro curto de uma decisão de arquitetura e de seus motivos \\
  \end{glossario}
\end{frame}

\begin{frame}{Questões de revisão}
  \begin{questoes}
    \item Que problemas aparecem quando uma \cod{Activity} faz tudo? Cite três.
    \item Compare MVP e MVVM: quem atualiza a tela em cada um?
    \item Por que o \cod{UiState} deve ser imutável e único por tela?
    \item Enuncie a regra da dependência e dê um exemplo de violação.
    \item Como a interface \cod{FonteLeituras} permite trocar BLE por MQTT sem mexer no \cod{ViewModel}?
    \item Qual a diferença entre injeção de dependência e localizador de serviços?
    \item Que escopo dar ao cliente HTTP, ao banco e a um \cod{ViewModel}? Por quê?
    \item Quando vale a pena criar a camada de domínio com casos de uso?
  \end{questoes}
\end{frame}

\begin{frame}{Exercícios}
  \begin{exercicio}[Exercício 1 --- Diagnóstico]
    Analise um projeto antigo seu (ou um exemplo da internet) e liste as responsabilidades misturadas. Desenhe como ficaria em camadas.
  \end{exercicio}
  \begin{exercicio}[Exercício 2 --- UiState]
    Modele o \cod{UiState} da tela de reserva: bancada escolhida, data, horário inicial e final, validação de conflito, envio em andamento e mensagem de erro.
  \end{exercicio}
  \begin{exercicio}[Exercício 3 --- Fake]
    Crie \cod{FakeBancadaRepository} que guarda os dados em um \cod{MutableStateFlow}. Use-o para alimentar o \cod{@Preview} e, no Módulo 14, os testes do ViewModel.
  \end{exercicio}
  \begin{projeto}[LabMonitor --- Etapa 8]
    Refatore o app para MVVM com Hilt: \cod{PainelViewModel}, \cod{EquipamentosViewModel}, repositórios com implementação \emph{fake} e o caso de uso \cod{VerificarAlertasUseCase}.
  \end{projeto}
\end{frame}

\begin{frame}{Para aprofundar}
  \small
  \begin{itemize}
    \item Guia de arquitetura de apps: \url{https://developer.android.com/topic/architecture}
    \item Recomendações de arquitetura: \url{https://developer.android.com/topic/architecture/recommendations}
    \item Camada de UI: \url{https://developer.android.com/topic/architecture/ui-layer}
    \item Camada de dados: \url{https://developer.android.com/topic/architecture/data-layer}
    \item Hilt: \url{https://developer.android.com/training/dependency-injection/hilt-android}
    \item Now in Android (app de referência): \url{https://github.com/android/nowinandroid}
  \end{itemize}
\end{frame}

\begin{frame}{Respostas das questões de revisão}
  \begin{respostas}
    \item A rotação refaz requisições e perde dados; as regras não podem ser testadas sem o app; uma mudança na API ou no banco quebra a tela; o código fica enorme e difícil de entender.
    \item No MVP, o \emph{presenter} chama métodos de uma interface da \emph{view}. No MVVM, o \cod{ViewModel} expõe estado e a \emph{view} o observa, sem que o \cod{ViewModel} a conheça.
    \item Um único objeto imutável evita estados contraditórios e atualizações pela metade; cada mudança gera uma cópia nova, fácil de comparar, registrar e testar.
    \item As camadas internas não conhecem as externas. Violação: um caso de uso que importa \cod{android.content.Context} ou classes do Retrofit.
    \item O \cod{ViewModel} depende só da interface, e o Hilt injeta \cod{FonteBle} ou \cod{FonteMqtt}: trocar é mudar uma linha no módulo de DI. Nos testes, entra uma \cod{FonteFake}.
    \item Na DI, a classe recebe as dependências de fora, explícitas no construtor. No localizador, ela mesma as pede a um registro global, o que as esconde e só falha em execução.
    \item Cliente HTTP e banco: únicos no app (\cod{@Singleton}), pois são caros e compartilhados. \cod{ViewModel}: ligado à tela, vive enquanto ela estiver na pilha.
    \item Quando regras se repetem em mais de um \cod{ViewModel}, são complexas ou precisam de testes isolados. Em apps pequenos, UI e dados bastam.
  \end{respostas}
\end{frame}

\encerramento{Módulo 09 (comum) --- Fundamentos de banco de dados: modelagem, SQL, normalização, transações e índices.}

\end{document}
